% TODO make hw stylesheet
\documentclass[10pt]{article}
\usepackage[letterpaper,top=1in,left=1in,right=1in]{geometry}
\usepackage[dvipsnames,svgnames]{xcolor}
\usepackage[shortlabels]{enumitem}
\usepackage[framemethod=TikZ]{mdframed}
\usepackage{amsmath,amssymb,amsthm}
\usepackage[colorlinks]{hyperref}
\usepackage{multicol}
\usepackage{tabularx}
\usepackage{bbm}

\hypersetup{
    colorlinks=true,
    linkcolor=blue,
    filecolor=magenta,
    urlcolor=magenta,
    pdftitle={MAT 549 HW}
}

\usepackage{mathtools}
\usepackage{thmtools}
\usepackage[textsize=footnotesize]{todonotes}
\usepackage{titling}
\usepackage{cancel}

\reversemarginpar
\mdfdefinestyle{mdbluebox}{roundcorner=10pt,innerbottommargin=9pt,
linecolor=blue,backgroundcolor=TealBlue!5,}
\declaretheoremstyle[headfont=\sffamily\bfseries\color{MidnightBlue},
mdframed={style=mdbluebox},]{thmbluebox}

\mdfdefinestyle{mdbloobox}{frametitlefont=\bfseries,innerbottommargin=8pt,
nobreak=true,backgroundcolor=TealBlue!5,linecolor=blue,}
\declaretheoremstyle[headfont=\bfseries\color{MidnightBlue},
mdframed={style=mdbloobox},headpunct={\\[3pt]},postheadspace=0pt,]{thmbloobox}

\mdfdefinestyle{mdredbox}{frametitlefont=\bfseries,innerbottommargin=8pt,
nobreak=true,backgroundcolor=Salmon!5,linecolor=RawSienna,}
\declaretheoremstyle[headfont=\bfseries\color{RawSienna},
mdframed={style=mdredbox},headpunct={\\[3pt]},postheadspace=0pt,]{thmredbox}

\mdfdefinestyle{mdgreenbox}{linecolor=ForestGreen,backgroundcolor=ForestGreen!5,
linewidth=2pt,rightline=false,leftline=true,topline=false,bottomline=false,}
\declaretheoremstyle[headfont=\bfseries\sffamily\color{ForestGreen!70!black},
mdframed={style=mdgreenbox},headpunct={ --- },]{thmgreenbox}

\mdfdefinestyle{mdorangebox}{linecolor=RedOrange,backgroundcolor=RedOrange!5,
linewidth=2pt,rightline=false,leftline=true,topline=false,bottomline=false,}
\declaretheoremstyle[headfont=\bfseries\sffamily\color{RedOrange!70!black},
mdframed={style=mdorangebox},headpunct={ --- },]{thmorangebox}

\mdfdefinestyle{mdblackbox}{linecolor=black,backgroundcolor=RedViolet!5!gray!5,
linewidth=3pt,rightline=false,leftline=true,topline=false,bottomline=false,}
\declaretheoremstyle[mdframed={style=mdblackbox}]{thmblackbox}

\declaretheoremstyle[spaceabove=6pt,spacebelow=6pt]{boldhead}
\declaretheoremstyle[spaceabove=6pt,spacebelow=6pt,headfont=\normalfont\itshape]{italichead}

\declaretheorem[style=boldhead,name=Problem]{problem}
\declaretheorem[style=boldhead,name=Problem,numbered=no]{problem*}
\declaretheorem[style=boldhead,name=Setup]{setup} % for config geo
\declaretheorem[style=boldhead,name=Example,sibling=problem]{example}
\declaretheorem[style=boldhead,name=$\bigstar$ Problem,sibling=problem]{niceprob}
\declaretheorem[style=boldhead,name=Theorem,sibling=problem]{theorem}
\declaretheorem[style=boldhead,name=Corollary,sibling=theorem]{corollary}
\declaretheorem[style=boldhead,name=Proposition,sibling=theorem]{prop}
\declaretheorem[style=boldhead,name=Lemma,numberwithin=problem]{lemma}
\declaretheorem[style=boldhead,name=Theorem,numbered=no]{theorem*}
\declaretheorem[style=boldhead,name=Lemma,numbered=no]{lemma*}
\declaretheorem[style=boldhead,name=Claim,numberwithin=problem]{claim}
\declaretheorem[style=boldhead,name=Claim,numbered=no]{claim*}
\declaretheorem[style=boldhead,name=Remark,numberwithin=problem]{remark}
\declaretheorem[style=boldhead,name=Remark,numbered=no]{remark*}
\declaretheorem[style=boldhead,name=Definition,sibling=theorem]{definition}
\declaretheorem[style=boldhead,name=Definition,numbered=no]{definition*}

\providecommand{\ol}{\overline}
\providecommand{\ceil}[1]{\left\lceil #1 \right\rceil}
\providecommand{\floor}[1]{\left\lfloor #1 \right\rfloor}
\newcommand{\dd}{\mathrm{d}}
\providecommand{\ul}{\underline}
\providecommand{\eps}{\varepsilon}
\providecommand{\half}{\frac{1}{2}}
\providecommand{\CC}{\mathbb C}
\providecommand{\FF}{\mathbb F}
\providecommand{\II}{\mathbb I}
\providecommand{\NN}{\mathbb N}
\providecommand{\QQ}{\mathbb Q}
\providecommand{\RR}{\mathbb R}
\providecommand{\ZZ}{\mathbb Z}
\providecommand{\ind}{\mathbbm 1}
\providecommand{\dg}{^\circ}
\providecommand{\ii}{\item}
\providecommand{\setdiff}{\smallsetminus}
\providecommand{\alert}[1]{{\sffamily\textbf{\textcolor{blue}{#1}}}}
\providecommand{\inv}{^{-1}}
\providecommand{\mb}[1]{\mathbf{#1}}
\newcommand{\what}{\widehat}

\newenvironment{soln}{\begin{proof}[Solution]}{\end{proof}}

\providecommand{\pow}{\operatorname{Pow}}
\providecommand{\vocab}[1]{{\textbf{\textcolor{ForestGreen}{#1}}}}
\providecommand{\lcm}{\operatorname{lcm}}
\providecommand{\abs}[1]{\left\lvert #1\right\rvert}
\providecommand{\smabs}[1]{\lvert #1\rvert}
\providecommand{\innerprod}[1]{\langle\!\langle #1\rangle\!\rangle}
\providecommand{\norm}[1]{\left\| #1\right\|}
\providecommand{\smnorm}[1]{\| #1\|}
\providecommand{\gr}{\operatorname{gr}}

\usepackage{mathrsfs}
\providecommand{\tanvec}[1]{\frac{\partial}{\partial #1}}

% place title at top right
\setlength{\droptitle}{-8ex}
\pretitle{\begin{flushright}\Large\bfseries}
\posttitle{\par\end{flushright}}
\preauthor{\begin{flushright}}
\postauthor{\end{flushright}}
\predate{\begin{flushright}}
\postdate{\end{flushright}}

\title{MAT 549 HW06}
\author{\large Aditya Pahuja}
\date{\large Due 04/21}
\begin{document}
\maketitle

\begin{problem}
    [16-3(a)]
    Suppose $E$ and $M$ are smooth $n$-manifolds
    with or without boundary, and $\pi\colon E\to M$ is
    a smooth $k$-sheeted covering map or generalized covering map.
    Show that if $E$ and $M$ are oriented and
    $\pi$ is orientation-preserving, then
    \begin{equation*}
        \int_E\pi^*\omega = k\int_M\omega
    \end{equation*}
    for any compactly supported $n$-form $\omega$ on $M$.
\end{problem}

\begin{soln}
    Let $\mathcal U$ be an open cover of $M$ by connected
    evenly covered open sets and let $\{\varphi_U : U\in\mathcal U\}$
    be a partition of unity subordinate to this open cover.
    Then, $\tilde{\mathcal U} = \{\tilde U = \pi\inv(U) : U\in\mathcal U\}$
    is an open cover of $E$ and the functions
    $\psi_{\tilde U} = \varphi_U\circ\pi$ for $\tilde U\in\tilde{\mathcal U}$
    form a partition of unity on $E$ subordinate to $\tilde{\mathcal U}$.
    To see that this collection of functions is indeed a partition of unity,
    we of course have that $\psi_{\tilde U}(x) = \varphi_U(\pi(x))\in[0,1]$
    and
    \begin{equation*}
	\sum_{\substack{\tilde U\in\tilde{\mathcal U}\\ x\in\tilde U}}
	\psi_{\tilde U}(x) =
	\sum_{\substack{U\in{\mathcal U}\\ \pi(x)\in U }}
	\varphi_U(\pi(x)) = 1
    \end{equation*}
    for all $x\in E$. Thus
    \begin{equation*}
        \int_E \pi^*\omega
	= \sum_{\tilde U\in\tilde{\mathcal U}}
	\int_{\tilde U}\psi_{\tilde U}\cdot\pi^*\omega
	= \sum_{\tilde U\in\tilde{\mathcal U}}
	\int_{U}\sum_{i=1}^k\varphi_{U}\cdot\omega
	= k\sum_{U\in\tilde{\mathcal U}}\int_U \varphi_U\cdot\omega
	= k\int_M\omega,
    \end{equation*}
    where the second equality is due to the fact that
    $\pi\inv(U)$ consists of $k$ components diffeomorphic to $U$ via
    the orientation-preserving map $\pi$
    and that $\psi_{\tilde U} = \varphi_U\circ\pi$.
\end{soln}

\begin{problem}
    [16-4]
    Suppose $M$ is an oriented compact smooth manifold with boundary.
    Show that there does not exist a smooth retraction
    of $M$ onto its boundary.
\end{problem}

\begin{soln}
    Suppose a smooth retraction $r\colon M\to\partial M$ exists.
    Since $M$ is orientable, so is its boundary,
    so let $\omega$ be an orientation form on $\partial M$,
    which in particular means that $\int_{\partial M}\omega \ne 0$.
    Then, Stokes's theorem gives the following contradiction:
    \begin{equation*}
	\int_{\partial M}\omega
	= \int_{\partial M} r^*\omega
	= \int_M \dd (r^*\omega) = \int_M r^*(\dd\omega) = 0.
    \end{equation*}
    To be thorough, the first equality holds because
    the pullback of $r^*\omega$ by the inclusion
    $\partial M\hookrightarrow M$ gives $\omega$
    because $r$ composed with the inclusion is the identity
    and the last equality is true because $\omega$ is a top-degree form
    on $\partial M$, hence $\dd\omega = 0$.
\end{soln}

\pagebreak
\begin{problem}
    [16-5]
    Suppose $M$ and $N$ are oriented, compact, and connected
    smooth manifolds, and $F,G\colon M\to N$ are homotopic diffeomorphisms.
    Show that $F$ and $G$ are either both orientation-preserving
    or both orientation-reversing.
\end{problem}

\begin{soln}
    Let $\omega$ be an orientation form on $N$, $I = [0,1]$,
    and let $H\colon M\times I\to N$ be a homotopy from $F$ to $G$,
    which we can take to be smooth by Theorem 6.29
    and which satisfies $H(p,0) = F(p)$ and $H(p,1) = G(p)$ for all $p\in M$.
    By Stokes's theorem and the fact that $\dd\omega = 0$,
    \begin{equation*}
	\int_{\partial(M\times I)} H^*\omega
	= \int_{M\times I} \dd H^*\omega
	= \int_{M\times I} H^*\dd\omega = 0.
    \end{equation*}
    Then, $\partial(M\times I) = M\times\{0,1\})$
    with the map $i\colon M\times\{0\}\to M\times\{1\}$
    given by $i(x,0) = (x,1)$ being orientation-reversing
    (due to the orientation that $\partial I = \{0,1\}$ inherits from $I$).
    Then, if $\iota_k \colon M\to M\times\{k\}$ is the map
    $\iota_k(x) = (x,k)$ for $k\in\{0,1\}$,
    \begin{align*}
	0 &= \int_{M\times I}H^*\omega
	= \int_{M\times\{0\}} H^*\omega - \int_{M\times\{1\}} H^*\dd\omega
	= \int_M \iota_0^* H^*\omega - \iota_1^* H^*\omega\\
	&= \int_M (H\circ\iota_0)^*\omega - (H\circ\iota_1)^*\omega
	= \int_M F^*\omega - G^*\omega,
    \end{align*}
    which implies that $F$ and $G$ either both preserve
    or both reverse orientation
    (since $\int_M F^*\omega = \int_M G^*\omega = \pm\int_N \omega$,
    with the sign depending on whether or not $F$ and $G$ reverse orientation).
\end{soln}

\begin{problem}
    [16-6]
    Prove that the following are equivalent.
    \begin{enumerate}[(a)]
        \ii There exists a nowhere-vanishing vector field on $S^n$.
	\ii There exists a continuous map $V\colon S^n\to S^n$
	satisfying $V(x)\perp x$ for all $x\in S^n$.
	\ii The antipodal map $\alpha\colon S^n\to S^n$
	is homotopic to the identity map.
	\ii The antipodal map $\alpha\colon S^n\to S^n$
	is orientation-preserving.
	\ii $n$ is odd.
    \end{enumerate}
\end{problem}

\begin{soln}
    (a) implies (b). This follows from the basic fact of Euclidean geometry
    that the unit normal at a given point $x\in S^n$ is $x$:
    if $\hat V$ is a nowhere zero vector field on $S^n$,
    then let $V\colon S^n\to S^n$ be defined by
    $V(x) = \frac{\hat V(x)}{\smnorm{\hat V(x)}}$,
    so that $V(x)$ is tangent to $S^n$ at $x$, hence orthogonal to $x$.\\

    (b) implies (c).
    Let $H\colon S^n\times I\to S^n$ be the map
    \begin{equation*}
        H(x,t) = \cos(\pi t)x + \sin(\pi t)V(x),
    \end{equation*}
    which is well-defined because
    $\norm{H(x,t)} = \cos(\pi t)^2 + \sin(\pi t)^2 = 1$ for all
    $(x,t)\in S^n\times I$
    and satisfies $H|_{S^n\times\{0\}} = \operatorname{Id}_{S^n}$,
    $H|_{S^n\times\{1\}} = \alpha$,
    so $H$ is a homotopy betwen $\alpha$ and the identity map.
    \\

    (c) implies (d). This follows from Problem 16-5
    and the fact that the identity map is obviously orientation-preserving.
    \\

    (d) implies (e). This follows from Problem 15-3
    from last week's homework.
    \\

    (e) implies (a). For $n = 2k - 1$ odd, let
    $V$ be the vector field on $S^n\subseteq\RR^{2k}$ defined by
    \begin{equation*}
        V(x_1,y_1,x_2,y_2,\dots,x_k,y_k) =
	(-y_1,x_1,-y_2,x_2,\dots,-y_k,x_k),
    \end{equation*}
    which is a tangent vector field because it
    satisfies $V(x)\perp x$ for all $x\in S^n$, and is nowhere zero
    because $\norm{V(x)} = \norm{x} = 1$ for all $x\in S^n$.
\end{soln}










\end{document}
